\section{Conclusion}
We presented A$^3$Bench, a benchmark of 2,683 videos and over 15K diagnostic questions designed to evaluate whether MLLMs can perform aggressive action and association -- recognizing not just that aggression occurs, but who initiates it, who is targeted, and what action connects them. Evaluation of 15 models reveals that current systems plateau below 50\% accuracy and degrade sharply as questions move from single-attribute recognition to full aggressor-action-victim reasoning. A consistent aggressor bias, in which models mislabel victims and bystanders as aggressors, persists across architectures and is not resolved by model scale or reasoning-style prompting. Our role-graph prompting strategy, which decomposes scene understanding into a structured causal chain of interaction detection, direction assignment, and role derivation, improves accuracy by up to 10 percentage points without any training, demonstrating that making the required reasoning decomposition explicit can partially compensate for the grounding deficits we identify. These results establish aggressive action and association as an open challenge for video-language models and provide a controlled diagnostic framework for measuring progress.
